\begin{recipe}{Quiche met geroosterde groente en geitenkaas}
\kicker[Hoofdgerecht,Vegetarisch,Provençaals]{Hoofdgerecht · vegetarisch · Provençaals}
\index[register]{Quiche met geroosterde groente en geitenkaas}
\index[register]{Geitenkaas}
\index[register]{Paprika}
\index[register]{Venkel}
\index[register]{Aubergine}

\meta{100 minuten \dvd Voor 6 porties \dvd Oven 180\,°C}

\lettrine{E}{en} quiche met geroosterde zomergroente, zachte geitenkaas en
provençaalse kruiden. Het deeg is hetzelfde als bij de andere quiches in dit boek.

\blockrule{Ingrediënten}

\noindent\textit{Deeg}
\begin{ingredients}
  \ing{300 g}{tarwebloem}
  \ing{150 g}{roomboter}
  \ing{2}{eieren}
  \ing{1 tl}{zout}
\end{ingredients}

\medskip
\noindent\textit{Vulling}
\begin{ingredients}
  \ing{75 g}{zachte geitenkaas}\index[register]{Geitenkaas}
  \ing{1}{rode paprika}\index[register]{Paprika}
  \ing{1}{gele paprika}
  \ing{1}{courgette}\index[register]{Courgette}
  \ing{1}{aubergine}\index[register]{Aubergine}
  \ing{1}{venkelknol}\index[register]{Venkel}
  \ing{2}{rode uien}
  \ing{2}{knoflookteentjes, geplet}
  \ing{200 g}{crème fraîche}
  \ing{4}{eieren}
  \ing{75 g}{geraspte kaas}
  \ing{100 g}{cashewnoten}
  \ing{1 takje}{rozemarijn}
  \ing{1 takje}{tijm}
  \ing{1 el}{Provençaalse kruiden}
  \ingb{zout}
\end{ingredients}

\blockrule{Bereiding}
\tip{Leg de groente niet te dicht op elkaar op de bakplaat, gebruik zo nodig
twee platen. Als de stukken elkaar raken gaan ze stomen in plaats van
roosteren en blijft de smaak vlak.}
\begin{steps}
  \item Kneed de deegingrediënten tot een samenhangend deeg. Maak een bol, wikkel
        in vershoudfolie en laat 15 minuten rusten in de koelkast.
  \item Verwarm de oven voor op 200\,°C (boven- en onderwarmte). 
  \item Was de paprika's, courgette, aubergine en venkel. Snijd ze samen met de rode ui in grove stukken.
        Leg op een bakplaat met de knoflook, rozemarijn en tijm.
  \item Schenk een scheutje olijfolie erover en rooster 25--30
        minuten tot de groente gaar en licht gekleurd is. Laat de groente
        daarna even uitlekken op de bakplaat: geroosterde courgette en
        aubergine geven nog vocht na.
  \item Draai de oven terug naar 180\,°C. Rol het deeg uit en bekleed de springvorm.
        Prik met een vork wat gaatjes in de bodem: dat voorkomt dat het deeg
        tijdens het voorbakken opbolt. Bedek de zijkanten eventueel met
        aluminiumfolie zodat ze niet naar beneden zakken. Bak de bodem
        15 minuten voor.
  \item Klop de eieren met de crème fraîche, de geraspte kaas en de provençaalse
        kruiden tot een glad mengsel. Breng flink op smaak met zout. Verkruimel
        de geitenkaas erdoor.
  \item Schep de geroosterde groente in de voorgebakken bodem. Giet het
        eimengsel erover.
  \item Bak 45 minuten op 180\,°C tot de vulling gestold is. Strooi de
        cashewnoten er na ongeveer 20 minuten over: zo roosteren ze mee zonder
        te verbranden.
  \item Laat de quiche 10 minuten rusten voor het aansnijden.
\end{steps}

\heroimagefade{images/quiche-groente-geitenkaas}
\end{recipe}
